\section{هفته اول}

\subsection{پنج‌شنبه ۱۷ آبان}
این اولین کامیتی هست که دارم روی این دفترچه خاطرات می‌زنم. امروز بالاخره دوتا از تسک‌های مهمی که برای خودم داشتم انجام شدن.
هم تونستم مقاله موردعلاقه خودم راجع به \lr{LUG} رو ترجمه کنم و هم یک سری ارائه با \lr{beamer} آماده کردم تا لاگ دانشگاه آزاد رو راه بندازم.

آخرین هدفی که داشتم این بود که باید کانتریبیوتر پروژه گو بشم. هدف عجیب‌غریب و حتی دست نیافتنی به‌نظر میاد. البته الان که دارم این متن رو می‌نویسم
اینطور هست. ممکنه چند وقت دیگه ممکن شده باشه و اگر شما دارید این دفترچه خاطرات رو می‌خونید یعنی این اتفاق افتاده!
پس دیگه میتونم پز بدم:) بهترین حس دنیا برای من این هست که بتونم داخل پروژه بزرگی مثل گولنگ مشارکت کنم و به‌عنوان توسعه دهنده فعال توش باشم.
اینکه این کار سخته یا آسون برام مهم نیست. امروز این کار شروع شده و باید اتفاق بیوفته.

داخل این مستند روز به روز خاطراتم رو برای رسیدن به این هدف خواهم نوشت و امیدوارم روزی بعد از من علاقه‌مندان از این نوشته‌ها استفاده کنن.
امید دارم خوندن قصه مسیری که دارم میرم بهتون کمک کنه. سعی می‌کنم تا جایی که میتونم با جزئیات باشه.

امروز نشستم فکر کنم که چطور می‌تونم مسیرم رو شروع کنم. اول به‌ذهنم رسید یه سر به ایده‌های قبلیم از مشارکت تو گو بندازم.
ایدم مربوط به بسته \lr{rpc} بود. گفتم خب چطور باید این کد رو خوند؟ چی میشه اگه بتونیم به‌صورت دیاگرام و گرافیکی ارتباط توابع یک کد رو داشته باشی؟
باید خیلی خوب بشه! برای همین گشتم دنبال ابزار و \lr{plantuml} و \lr{goplantuml} رو پیدا کردم. اینطوری شد که میتونم کدی که دارم رو
خیلی راحت تر تحلیل کنم و با داشتن یک \lr{big image} بشینم باقی کد رو بخونم.
نتیجه نهایی من شد کد زیر:
\begin{latin}
\begin{minted}{bash}
goplantuml . | plantuml -darkmode -pipe | feh -
\end{minted}
\end{latin}

خب این کد خیلی خفنه واقعا. میری داخل پوشه ای که دوست داری و اجراش می‌کنی، بعد یک عکس بهت میده از تحلیل روابط توابع و کلاس‌های اون
کد گو. حتما پاشید برید راجع به \lr{plantuml} مطالعه کنید که بدردتون میخوره. نمیخوام وقتمو با گذاشتن عکس تلف کنم وگرنه میذاشتم ببینید.
ولی خب شما هم باید یکم سرچ کنید دیگه نمیشه که همش من سرچ کنم.

شبو رفتم بیرون و دیگه کار رو ادامه ندادم.

\subsection{جمعه ۱۸ آبان}

یکمی دیر شروع کردم امروز. ساعت ۱۶:۳۱ دقیقه نشستم پا کار و قراره مستند مشارکت خود گو رو بخونم و ستاپم رو اوکی کنم.
\LTRfootnote{\url{https://go.dev/doc/contribute}}

مراحل آماده کردن سیستم به‌عنوان یک توسعه‌دهنده گو خیلی جالبه. اول باید یک سری قرارداد رو بخونی و امضا کنی.
این قرارداد می‌گه مثلا ما میتونیم مشارکت‌هات رو عمومی کنیم و از این حرف‌ها. بعد از اینکه با ایمیل خودت قرارداد یا 
\ftnt{CLA}{https://golang.org/doc/contribute.html#cla}
رو امضا کردی میری سراغ مراحل بعدی. حالا باید بری و تو گیت سرور مخصوص پروژه گو که دست گوگل هست یه سری کوکی ثبت کنی و
گیتت رو کانفیگ کنی. وقتی میری داخل رابط گرافیکی گیت سرور
\LTRfootnote{\url{https://go.googlesource.com/}}
اونجا باید به‌حساب یه پسورد درست کنی و بعد یه اسکریپت مثل چیزی که این زیر هست بهت میده تا اجراش کنی:

\begin{latin}
\begin{minted}{bash}
eval 'set +o history' 2>/dev/null || setopt HIST_IGNORE_SPACE 2>/dev/null
touch ~/.gitcookies
chmod 0600 ~/.gitcookies

git config --global http.cookiefile ~/.gitcookies

tr , \\t <<\__END__ >>~/.gitcookies
SECRET TOKENS THAT I HID
__END__
eval 'set -o history' 2>/dev/null || unsetopt HIST_IGNORE_SPACE 2>/dev/null
\end{minted}
\end{latin}

وقتی اینکارو کنی به‌حساب کانفیگ کردی سیستمت رو. حالا باید بریم سراغ پلتفرم \lr{Gerrit}. یه پلتفرم کد ریویو هست خیلی هم جالبه.
به‌جای اینکه پچ هامون رو بفرستیم به یه لیست پستی، می‌فرستیم به \lr{gerrit} و اونجا \lr{reviewer} ها بررسی می‌کنن و کامنت میذارن و
از این جور کارها. خلاصه که اونجا هم باید اکانت بسازی حتما.

حالا واسه اینکه بتونیم پچ بفرستیم به \lr{gerrit} از یه دستور که باید پایین تر نصبش کنیم استفاده می‌کنیم به اسم \lr{git-codereview}

\begin{latin}
\begin{minted}{bash}
go install golang.org/x/review/git-codereview@latest
\end{minted}
\end{latin}

اینکارو که بکنید بعد با خود گیت می‌تونید استفاده کنید ازش. یه زیر دستور برای گیت میشه. چه‌جالب نمیدونستم از این کارا با گیت میشه کرد...
باید یه روز گیت رو هم خوب یاد بگیرم :))

مستند رو کامل خوندم. خب نیاز هست موردی بهش رجوع کنم ولی خب خیلی کامل و خفنه. پروسه مشارکت در گو واقعا عجیب غریبه. حداقل الان
اینطور فکر می‌کنم و ممکنه بعدا نظرم عوض بشه. اگر \lr{pdf} داشته باشه مستندش ۱۷ صفحه رو میشه.

زمانی که میخواید یه کار مهم انجام بدید بعد از اینکه ایشو زدید باید پروپوزال هم ارائه بدید. این فرایند هم طولانی هست. این تیپیه که مثلا
یک سری آدم هفته ای یک بار دور هم میشینن و بعد پروپوزال ها رو بررسی می‌کنن. در کل گیج شدم. احساس میکنم مشارکت تو پروژه گو
وحشتناک سخته. حتی حس می‌کنم از مشارکت تو لینوکس هم سخت تره و نمیشه به راحتی تو این پروژه مشارکت متن‌باز داشت.
ولی قول میدم بتونم این کارو انجام بدم. حتی داشتم فکر می‌کردم که چطوره به‌جای لاتک این دفترچه خاطرات رو داخل یک \lr{static site generator}
منتقل کنم و هر روز پست درست کنم. به‌نظر کار بهتری میاد و بیشتر استقبال می‌شه. فعلا همینجا کارم رو ادامه میدم و هر زمان شد اون کار رو هم
انجام میدم حتما. اول باید بشینم یه چیزی مثل \lr{Hugo} یاد بگیرم.

خلاصه میشه که گیجم گیجم گیجم گیجم گیجم گیجم. به‌نظر واقعا غیر ممکن میاد. نمیشه همچین کاری انجام داد:)) میشه واقعا؟

\subsection{چهارشنبه ۱۹ دی}
دقیقا ۲ ماه و ۱ روز از اخرین فعالیتم راجع به این مسئله می‌گذره. ۲ ماه سخت رو درگیر امتحانات میان ترم و پایان ترم بودم ولی به هر نحوی گذشت.
الان دیگه درسی ندارم و میتونم برسم به کارم. البته مشکل اینجاست که مجددا باید برم مستندات و \lr{aggrement} ها رو مطالعه کنم.
ایرادی نداره، اینکار رو انجام می‌دم و دوباره میخونمشون و مسیرم رو مجددا شروع خواهم کرد.

اول از همه یک نگاهی به \lr{Gerrit} میندازم و مستنداتش
\LTRfootnote{\url{https://gerrit-review.googlesource.com/Documentation/}}
رو مطالعه می‌کنم. فکر می‌کنم بخش های اول مستندش فعلا کافی باشه. بیشتر رو یا در آینده مطالعه می‌کنم یا هم که یاد می‌گیرم در حین کار.
ولی جالب هست آدم بدونه این سرویس مستقل هست و خیلی ها استفاده می‌کنن ازش. حالا پروژه گو هم یکی داره ازش و هرکسی میتونه
روی یونیکس خودش یک نمونه اجرا کنه و استفاده کنه ازش.

مستندات رو خوندم دوباره. حالا باید برای خودم چند تسک تعریف کنم و اونها رو انجام بدم. بعد از اون بشینم پای مشارکت.
ولی قبل از هرچیز باید با روند پروژه بیشتر آشنا بشم.

بخش های زیادی برای بررسی وجود داره. تسک ها رو لیست می‌کنم:
\begin{itemize}
\item
بررسی تمام زیرپروژه های
\LTRfootnote{\url{https://go.googlesource.com/}}
و تهیه خلاصه از عملکردشون.
\item
بررسی بیلد ها در سرور بیلد خود گو
\LTRfootnote{\url{https://build.golang.org/}}.
\item
بررسی انواع \lr{CL} ها با وضعیت‌های مختلف باز بسته و رها شده از داخل \lr{Gerrit}
\LTRfootnote{\url{https://go-review.googlesource.com/q/status:open+-is:wip}}.
\end{itemize}

از بین این موارد کامیت خوندن و بررسی تغییرات دیگران و کامنت ها خیلی خوب هست. این مورد رو باید خیلی انجام بدم.
حتی میتونه این یک تمرین و تفریح روزانه باشه که آدم خبردار باشه چه اتفاق هایی داره میوفته. البته که در شروع تسک سختی هست.

فعلا با کلون کردن و بررسی زیرپروژه های مختفل شروع می‌کنم. همینجا میگم هرکدوم چی هستن و چیکار میکنن.
از اسکریپت زیر برای کلون کردن همه زیر پروژه ها و درست کردن یه todo list واسه خودم استفاده می‌کنم. یک دفعه یاد کسایی افتادم
که یک عالمه فیلم دانلود می‌کنن و نمی‌بینن:)) البته قراره من تمام این پروژه ها رو تا پایان امروز حداکثر بررسی کنم.

\begin{latin}
\begin{minted}{bash}
echo "You should review all of this projects" >> tasks.md

BASE=https://go.googlesource.com

for project in arch benchmarks blog build cbro-scratch crypto \
                debug dl example exp gddo go gofrontend gollvm \
                govulncheck-action grpc-review image lint mobile \
                mod net oauth2 open2opaque oscar perf pkgsite \
                pkgsite-metrics playground proposal protobuf \
                review scratch sublime-build sublime-config \
                sync sys talks telemetry term text time tools \
                tour vgo vscode-go vuln vulndb website wiki \
                xerrors
do
        git clone "$BASE/$project"
        echo "- [ ] $BASE/$project" >> tasks.md
done
\end{minted}
\end{latin}
بعد از اجرای این کد، حالا یه عالمه کد برای بررسی داریم. میشینم همشون رو طبق لیست بررسی می‌کنم. هرچقدر بتونم اینجا توضیح میدم.
هم برای یادگیری خودم خوبه هم در آینده میشه ازش استفاده کرد.

خیلی وقت ها برای اینکه بفهمم یه کد چیکار میکنه یا مثلا یک تابع چیکار میکنه، اسمش رو داخل پروژه های دیگه سرچ می‌کنم.
از این دستور استنفاده می‌کنم:

\begin{latin}
\begin{minted}{bash}
# find . -type f -name "*.go" -exec grep -nH "golang.org/x/arch" {} \;
./src/cmd/internal/disasm/disasm.go:28:	"golang.org/x/arch/arm/armasm"
./src/cmd/internal/disasm/disasm.go:29:	"golang.org/x/arch/arm64/arm64asm"
./src/cmd/internal/disasm/disasm.go:30:	"golang.org/x/arch/loong64/loong64asm"
./src/cmd/internal/disasm/disasm.go:31:	"golang.org/x/arch/ppc64/ppc64asm"
./src/cmd/internal/disasm/disasm.go:32:	"golang.org/x/arch/riscv64/riscv64asm"
./src/cmd/internal/disasm/disasm.go:33:	"golang.org/x/arch/s390x/s390xasm"
./src/cmd/internal/disasm/disasm.go:34:	"golang.org/x/arch/x86/x86asm"
\end{minted}
\end{latin}
مثلا اینجا میخواستم ببینم دقیقا این پروژه \lr{arch} چیکار میکنه؟ یه سری جدول میسازه از دستورات اسمبلی و مشخصاتشون تو معماری‌های
مختلف... خب به چه درد میخوره؟ وقتی این سرچ رو زدم خب متوجه شدم کجا باید بگردم دنبال جوابم. 
واسه همین میتونم استنتاج کنم این پکیج درواقع یکی از جاهایی که استفاده شده اونجایی هست که خط فرمان گو میاد دیس‌اسمبلی می‌کنه باینری‌ها رو.
استفادش هم این هست که اطلاعات راجع به سخت افزار های مختلف رو لیست کرده و میشه یک سری کار رو مختص به سخت افزار انجام داد.

من واسه اینکه ببینم مثلا چیکار میشه کرد با توابعش یه کد کوچیک زدم. خیلی جالبه. دستورات اسمبلی جهت یک سری استفاده که خودم هنوز نمیدونم
انکد میشن و بعد میتونیم با توابعی که اینجا آماده شده مثلا اینا رو دیکد کنیم. نمونه کدم رو این زیر میذارم. اینجا یک سری کد هگز داریم و مود پردازنده.
بعد میتونیم سینتکس های \lr{GNU}, \lr{Intel} و \lr{Go} رو با این کد دید. خود اسمبلی سینتکس اسمبلی خودش رو داره. خیلی جالبه که سینتکس اسمبلی
گو مشابه اسمبلی هست که در \lr{Plan 9} سیستم عامل قدیمی که خود آقای کن تامسون از توسعه دهنده های زبان سی و الان زبان گو بوده.
مثل اینکه کن تامسون جان هنوز هم به پروژه های قدیمیش علاقه داره و دوست نداره فراموش بشن. اینکه سینتکس اسمبلی رو مثلا گنو میذاشتن خیلی
سازگاری بیشتری داشت و من هم حداقل آشنا بودم باهاش... ایرادی نداره. داخل مستندات \lr{plan9} سینتکس اسمبلیش هم توضیح داده شده.
\LTRfootnote{\url{https://9p.io/sys/doc/asm.html}}

\begin{latin}
\begin{minted}{go}
package main

import (
        "encoding/hex"
        "fmt"

        "golang.org/x/arch/x86/x86asm"
)

func main() {
    for _, code := range []string{
        "556677885f5f5f5f5f5f",
        "223344556677885f5f5f5f5f5f",
        "223344556677885f5f5f5f5f5f",
        "11223344556677885f5f5f5f5f5f",
        "223344556677885f5f5f5f5f",
    } {
        for _, mode := range []int{64, 32} {
            c, _ := hex.DecodeString(code)
            ins, _ := x86asm.Decode(c, mode)

            fmt.Printf("\033[31mgo    (%d): %s\n", mode, x86asm.GoSyntax(ins, 0, nil))
            fmt.Printf("\033[33mGNU   (%d): %s\n", mode, x86asm.GNUSyntax(ins, 0, nil))
            fmt.Printf("\033[36mIntel (%d): %s\n", mode, x86asm.IntelSyntax(ins, 0, nil))
        }
    }
}
\end{minted}
\end{latin}

به این ترتیب بررسی پکیج \lr{arch} تموم میشه. چیزای جالبی ازش یاد گرفتم. تازه دوتا موضوع برای مشارکت هم به ذهنم رسید.
ساختار \lr{Makefile} اش نامرتب بود. فکر کنم بشه تغییرش داد و فرستاد. اینکار رو شده برای تمرین انجام میدم.
تازه اضافه کردن یک فایل \lr{.gitignore} هم به‌نظرم بد نمیاد. کدبیس قدیمی هست. تازه همین الان دیدم یک دوست عزیزی
پاشده غلط املایی گرفته از داخل کامنت دقیقا تو همین پروژه:)))
\LTRfootnote{\url{https://go.googlesource.com/arch/+/bde81be39b9efbf5b80719bb91ad1a0ebc5186b5}}
حتی همین هم بشه من انجام میدم. همنقدر الکی میشه مشارکت کرد واقعا. حتی بررسی کد برای اینطور چیز ها هم ایده بدی نیست.
بهرحال شناخت \lr{workflow} توسعه متن‌باز از یکجایی باید شروع بشه. حداقل با اینکار میشه یک سری چیزای کوچیک رو فهمید.

درکل یادم باید باشه که پکیج \lr{arch} کدبیس قدیمی داره و خیلی میشه بهترش کرد. کلا کم روش کار شده. یه کدبیس زدن همه
هم میترسم برن دست بزنن بهش:) همین خوبه واسه شروع کاری. میشه روش فکر کرد.

\subsection{چهارشنبه ۲۰ دی}
امروز خیلی تنبلی کردم و دیر دارم کار رو شروع میکنم. در عوض شب با دوستام بیرون نمیرم و قول میدم زود بخوابم.
فضای مجازی هم تا فردا کلا کنسله. بهرحال اگر نشه خودم رو تنبیه کنم این زندگی درست بشو نیست. حتی برای دانلود موزیک هم کنسله...

بگذریم بریم سر کار و بررسی پروژه \lr{benchmark}. دیروز فکر میکردم بتونم خیلی بیشتر از این حرفا پروژه بررسی کنم ولی فقط تونستم \lr{arch}
رو بررسی کنم. ۵۰ تا پروژه هست و اکثرا هم پیچیده هستن. خدا میدونه چقدر تایم میگیره... ولی مجبورم از یک جا به بعد از سر وا کنم.
نباید خیلی سر شناخت اینا تایم بذارم. شاید شکوندمشون بعدا و در آینده بیشتر بررسی کردمشون.

با بررسی این پروژه متوجه شدم که ساختار بنچمارکینگ گو از این پکیج استفاده می‌کنه. درواقع \lr{entrypoint} های ما موقع تست کارایی اینجا تعریف
میشن. بماند که یک سری کد برای تست عملکرد خود کامپایلر هم وجود داره. تمام این پروژه ها مربوط به ساختار های داخلی گو هستن.
برای مثال همین پروژه موضوع مهمش مربوط به درایورش هست که یک مثال ساده سعی میکنم باهاش پیاده کنم و جلو تر بیارم.

\begin{latin}
\begin{minted}{go}
package main

import (
  "encoding/hex"

  "golang.org/x/benchmarks/driver"
)

func main() {
  driver.Main("HEX", bechmark)
}

func bechmark() driver.Result {
  return driver.Benchmark(benchmarkN)
}

func benchmarkN(N uint64) {
  driver.Parallel(N, 4, func() {
    _, err := hex.DecodeString("a123b34b23bc23d233f0")
    if err != nil {
      panic(err)
    }
  })
}
\end{minted}
\end{latin}

خروجی این کد به صورت زیر است:
\begin{latin}

\begin{minted}[fontsize=\footnotesize,breaklines]{text}
pkg: golang.org/x/benchmarks
goos: linux
goarch: amd64

2025/01/09 16:07:56 Benchmarking 1 iterations
2025/01/09 16:07:56 Benchmarking 100 iterations
2025/01/09 16:07:56 Benchmarking 10000 iterations
2025/01/09 16:07:56 Benchmarking 1000000 iterations
2025/01/09 16:07:57 Benchmarking 100000000 iterations
2025/01/09 16:08:01 Benchmarking 200000000 iterations
# memprof=/tmp/20.prof.txt
# cpuprof=/tmp/19.prof.txt
BenchmarkHEX-2 200000000                43 ns/op           2894096 GC-bytes-from-system     234595 STW-ns/GC             0 STW-ns/op            16 allocated-bytes/op            1 allocs/op      47273224 bytes-from-system      41435136 heap-bytes-from-system          2485240 other-bytes-from-system        36765696 peak-RSS-bytes        1258930176 peak-VM-bytes            458752 stack-bytes-from-system              73 user+sys-ns/op
\end{minted}
\end{latin}

به همین ترتیب عملکرد \lr{benchmark} هم به‌عنوان زیرساخت تست کارایی در گو مشخص می‌شه. بنچمارک هامون رو هم میتونیم
این شکلی مستقیم با \lr{entrypoint} بنویسیم. ولی خب خود گو اینکار رو برامون انجام میده.

